% ECONOMICS_PROBLEM: {"id": "C78-424", "title": "Efficient Birkhoff representation of two-sided robust stable matchings", "jel_code": "C78", "source_id": "OP-0198"}
% !TeX program = xelatex
\documentclass[11pt,a4paper]{article}
\usepackage[margin=23mm]{geometry}
\usepackage{fontspec}
\setmainfont{Latin Modern Roman}
\usepackage{amsmath,amssymb,mathtools}
\usepackage{xcolor,enumitem,booktabs,longtable,array,xurl,hyperref}
\definecolor{muted}{HTML}{53636C}
\hypersetup{colorlinks=true,urlcolor=blue,linkcolor=blue}
\setlength{\parindent}{0pt}
\setlength{\parskip}{5pt}
\newcommand{\R}{\mathbb R}
\newcommand{\E}{\mathbb E}
\newcommand{\N}{\mathbb N}
\newcommand{\ind}{\mathbf 1}
\newcommand{\Var}{\operatorname{Var}}
\newcommand{\Cov}{\operatorname{Cov}}
\newcommand{\supp}{\operatorname{supp}}
\DeclareMathOperator*{\argmin}{arg\,min}
\DeclareMathOperator*{\argmax}{arg\,max}
\newcommand{\entrymeta}[3]{{\sffamily\small\color{muted}\textbf{\texttt{#1}}\quad|\quad #2\par #3\par}}
\newcommand{\Problemref}[1]{\texttt{#1}}
\newcommand{\JELref}[2]{\texttt{#2} (JEL \texttt{#1})}
\begin{document}
\section*{Efficient Birkhoff representation of two-sided robust stable matchings}
\textbf{Problem ID:} \texttt{C78-424}\quad \textbf{JEL:} \texttt{C78}\par
% BEGIN ECONOMICS STATEMENT
\label{problem:OP-0198}
% GAME_THEORY_PROBLEM: {"local_id": "GT-068", "original_number": 68, "kind": "P", "entry_kind": "statement_only", "published": false, "review_required": true, "category": "game_theory"}
\label{game:GT-068}
{\sffamily\small\color{muted}
\noindent\textbf{ID: GT-068.}\quad\textbf{Category:} Stable matching algorithms.
\quad\textbf{Kind: P.}\quad\textbf{Status:} Open in the cited source.
\par}
\subsection*{Definitions and mathematical statement}
Let \(W,F\) be disjoint sets of \(n\geq1\) workers and \(n\) firms.
A preference profile assigns each participant a strict total order on
the opposite set. A perfect matching is a bijection \(\mu:W\to F\).
It is stable for profile \(A\) if no \((w,f)\in W\times F\) satisfies
\[
 f\succ^A_w\mu(w)
 \quad\text{and}\quad
 w\succ^A_f\mu^{-1}(f).
\]
Let \(\mathcal M_A\) be the stable matchings for \(A\).
Suppose profiles \(A,B\) differ in at most one worker's list and at most
one firm's list, and put \(\mathcal R=\mathcal M_A\cap\mathcal M_B\).

Order \(\mathcal M_A\) by \(\mu\leq_A\nu\) if every worker weakly
prefers \(\mu(w)\) to \(\nu(w)\) in \(A\). The meet gives every worker
their preferred partner among the two matchings, and the join their less
preferred partner. When \(\mathcal R\ne\varnothing\), it is a finite
distributive sublattice. For a poset \(P\), an ideal is a set
\(D\subseteq P\) with \(x\in D,\ y\leq_P x\Rightarrow y\in D\).
The ideals form a lattice under intersection and union.

\paragraph{Question.}
Is there a polynomial-time algorithm that either reports
\(\mathcal R=\varnothing\), or outputs a polynomial-size poset \(P\)
and an effective lattice isomorphism between its ideals and
\((\mathcal R,\leq_A)\)? Effectiveness requires polynomial-time
conversion in both directions between an ideal and its matching,
with meet and join corresponding to intersection and union.

\subsection*{Short English statement}
Can all jointly stable matchings be represented efficiently by a poset?

\subsection*{Variants and scope boundaries}
Finding one jointly stable matching is already polynomial-time solvable
in this regime. Birkhoff's existence theorem alone does not supply the
required efficient representation. Explicitly listing all robust
matchings may require exponential space.

\subsection*{Source relationship and status}
[S1, Remark 16] asks for this Birkhoff partial order. The conversion
requirement above makes the algorithmic representation precise.
The catalogue's second author name is corrected to Tung Mai.

\subsection*{References}
\begingroup\footnotesize\raggedright
\noindent[S1] Rohith Reddy Gangam, Tung Mai, Nitya Raju, and
Vijay V.~Vazirani,
\emph{Stable Matching: Dealing with Changes in Preferences}, 2023;
arXiv:2304.02590v3, 2025.
\url{https://arxiv.org/html/2304.02590v3}.
Locators: Sections 3.1--3.2; Section 4.1, Corollary 1;
Section 4.4, Remark 16.
\par\endgroup

\subsection*{Common conventions: Source mathematical conventions}

\textbf{Finite games.} For a finite set $A$, $\Delta(A)$ is its probability
simplex. Players choose independent mixed strategies in Nash equilibrium;
correlation is allowed only when explicitly part of the solution concept.
In a game with payoff functions $u_i$ and strategy profile $x$, an additive
$\varepsilon$ Nash equilibrium satisfies
\[
 u_i(x)\geq u_i(x_i',x_{-i})-\varepsilon
 \quad\text{for every player $i$ and every admissible deviation $x_i'$}.
\]
For numerical approximation, payoffs are normalized as locally specified.
An $\varepsilon$ payoff residual is distinct from distance at most
$\varepsilon$ to the set of exact equilibria. Point convergence, convergence
of empirical play, and convergence in distance to an equilibrium set are
also distinct.

\textbf{Computational representations.} Unless stated otherwise, a complexity
question takes rational data with binary encoding; polynomial time is measured
in total bit length. A fixed-accuracy polynomial algorithm is not necessarily
polynomial in $1/\varepsilon$, and neither is necessarily polynomial in
$\log(1/\varepsilon)$. Succinct games and value, demand, or sampling oracles
must declare their interfaces. Exact real solutions require an explicit
representation; an unspecified list of arbitrary real numbers is not a
finite computational output. A claimed algorithmic barrier is meaningful
only for its specified model and reductions.

\textbf{Dynamic games.} Histories, information, observation, and allowable
randomization are part of the model. A stationary strategy depends only on
the current state; a positional strategy in a deterministic graph game
chooses one action at each controlled vertex. Nash equilibrium at an initial
state and equilibrium after every history are different requirements.
An expectation of a pathwise $\liminf$ payoff, a $\liminf$ of expected averages,
a limiting expected average, and a uniform finite-horizon equilibrium are
different criteria. None is silently substituted for another.

\textbf{Allocation and mechanisms.} A complete allocation assigns every item
exactly once unless otherwise stated. Zero-valued goods, negative values,
capacity constraints, feasibility, lotteries, and copies of goods affect
fairness definitions and are specified locally. Dominant-strategy incentive
compatibility, Bayesian incentive compatibility, universal truthfulness,
and truthfulness in expectation impose different inequalities. Whenever
an objective ratio has zero denominator, its convention is stated or those
instances are excluded.

\textbf{Infinite-dimensional models.} Expectations, integrals, stochastic
equations, and optimization objectives require their local regularity,
integrability, filtrations, and admissible domains. A backward--forward
mean-field system is a boundary-value problem; it is not an initial-value
semiflow without an additional construction. Long-time, large-population,
and vanishing-noise limits require an order of limits and a convergence
topology. If a minimum need not be attained, the objective is its infimum
or supremum and the approximation requirement is stated separately.
% END ECONOMICS STATEMENT
\end{document}
